\chapter{Conclusion}

\section{Summary of Context and Objectives}

This project set out to address a genuine gap in the landscape of
technical education tools: the absence of a platform capable of
generating structured, enriched programming quizzes and coding
challenges on demand, tailored to any topic a learner wishes to
practise. Existing solutions either offer static question banks that
grow stale, or expose raw AI generated content without any
organisational structure or quality signal. The platform developed in
this dissertation bridges that gap by combining the generative power
of large language models \cite{brown2020gpt3,touvron2023llama} with
two custom trained machine learning classifiers that automatically
enrich every generated question with a predicted topic tag and a
predicted difficulty level.

Six objectives were defined at the outset of the project. These
covered the full stack web application and its AI generation pipeline,
the two classifiers and their integration into the backend inference
layer, the multi language code execution engine, and the user account
system with its gamification and social features. All six objectives
were pursued simultaneously across an iterative development process
grounded in agile principles \cite{agilemanifesto2001,beck2004xp},
with a GitHub based workflow, continuous integration pipeline
\cite{githubactions2025}, and containerised Docker Compose
infrastructure \cite{merkel2014docker,dockercompose2025} supporting
the full development lifecycle. Each objective maps to a concrete
subsystem in the delivered software, and each has been evaluated
against measurable criteria in Chapter 5.

\section{Summary of Findings}

The system evaluation chapter assessed each objective against
measurable criteria. The key findings are summarised below.

The AI generation pipeline functioned reliably across both supported
backends. The OpenAI \texttt{gpt-4o-mini} backend consistently
produced clean, schema conforming JSON for both MCQ and coding
challenge quiz types. The Ollama \texttt{llama3.1:8b} backend
required additional post processing to handle formatting
inconsistencies in its output but performed correctly once these were
addressed. JSON parsing fragility when requesting larger numbers of
questions was identified as the primary limitation of this component,
and the structured output mode now offered by the OpenAI API
\cite{openai2024structured} offers a clear path to resolving it.

The topic tag classifier achieved a test set accuracy of 0.805 and an
F1 macro of 0.622 across ten algorithmic topic classes. Performance
was strong on classes with distinctive vocabulary, with Database
achieving a perfect F1 of 1.00 and Array achieving 0.91. Performance
was weaker on conceptually overlapping classes such as Hash Table
(0.20) and Dynamic Programming (0.35), where genuine ambiguity in the
label definitions limits the ceiling of any text based classifier
\cite{zhang2014multilabel}. The shrinkage regularised per label
threshold optimisation, informed by the approach of Lipton et al.
\cite{lipton2014thresholding}, produced noticeable gains over the
default 0.5 threshold on the minority classes.

The difficulty classifier achieved a test set accuracy of approximately
0.593 on a balanced three class problem. The XGBoost model
\cite{chen2016xgboost} outperformed all six alternatives on validation
accuracy. The medium difficulty class was the hardest to classify
across all models, which is consistent with its position as the
boundary class between easy and hard. The fundamental limitation of
this component is the proxy label: acceptance rate is a reasonable
but imperfect stand in for human assigned difficulty
\cite{zhou2017weaksupervision}.

Both classifiers were successfully integrated into the backend
inference pipeline as in process singletons, adding negligible latency
to the question generation flow. Every generated question is
automatically enriched with topic tags and a difficulty label before
being returned to the frontend, delivering on the fourth objective
directly.

The code execution engine supported all five target languages (Python,
JavaScript, TypeScript, Java, and C\#) with a five second timeout and
per test case pass or fail evaluation. Execution was reliable across
all languages, with Java and C\# requiring careful handling of runtime
entry point conventions.

The user accounts system, including JWT authentication \cite{jwt2015},
bcrypt password hashing \cite{provos1999bcrypt}, the experience and
levelling mechanic, and the full social friends lifecycle, functioned
correctly in all manual testing scenarios.

\section{Unexpected Insights}

Several insights emerged during the project that were not anticipated
at the outset.

The most significant was the discovery of how strongly class imbalance
dominates multi class text classification performance. It was initially
assumed that accuracy would be the primary indicator of classifier
quality, and it was only through the process of computing F1 macro
alongside accuracy that the disparity became apparent: a model scoring
0.805 accuracy was simultaneously scoring only 0.622 on F1 macro,
because it was essentially ignoring minority classes. This led to the
adoption of per label threshold tuning and the use of F1 macro as the
primary model selection criterion, which meaningfully changed which
model was chosen for deployment. The lesson that accuracy can conceal
minority class failure is well documented in the literature
\cite{japkowicz2002imbalance}, but encountering it in a concrete
project rather than in the abstract made the point land far more
forcefully.

A second unexpected insight came from the difficulty classifier's
label engineering. The decision to use acceptance rate as a proxy for
difficulty initially felt like a pragmatic workaround, but the process
of investigating why the medium class was consistently the hardest to
classify revealed something genuinely interesting about the nature of
difficulty itself: the boundaries between easy, medium, and hard are
not crisp in any dataset, human or otherwise. Problems cluster
continuously along a difficulty spectrum, and any binning strategy
introduces arbitrary boundaries. This is not unique to this project
and reflects a broader challenge in any ordinal classification problem
derived from continuous data.

A third observation, more practical in nature, was the extent to which
prompt engineering turned out to be as important as model selection for
the quality of the AI generation pipeline. Minor changes to the
structure and specificity of the prompt templates had a larger effect
on output consistency than switching between model backends. This
echoes findings from the broader prompt engineering literature
\cite{liu2023promptsurvey,wei2022chainofthought}, and it has concrete
implications for how the project's templates should be versioned and
tested going forward. The prompts are basically part of the product itself, so they should be handled like code, properly reviewed, version controlled, and tested to make sure they still give the expected results.

\section{Opportunities for Future Investigation}

Several directions for extending and improving the platform were
identified during the course of the project.

The most impactful improvement to the topic tag classifier would be
to expand the training dataset to include more examples of the
minority classes, or to acquire a dataset with a more balanced
distribution across tags. Oversampling techniques such as SMOTE
\cite{chawla2002smote} applied to the TF IDF feature space, or the
use of class weighted loss functions more aggressive than the current
\texttt{balanced} setting, could also be explored. The adoption of a
pre trained language model such as a fine tuned BERT variant
\cite{devlin2019bert} for feature extraction, rather than a bag of
words TF IDF approach, would likely yield stronger performance on
the conceptually overlapping classes where surface vocabulary is
insufficient to discriminate between labels.

For the difficulty classifier, the most meaningful improvement would
be the introduction of human verified difficulty labels, even for a
small subset of problems, to serve as ground truth for evaluating and
calibrating the proxy label approach. The current model could also be
extended to incorporate solution level features such as the complexity
of accepted solutions or the distribution of approaches used, which
would provide a richer signal than problem description text alone.
Replacing the three way classification with a regression onto a
continuous difficulty score, and binning only at presentation time,
would also remove the arbitrary boundary artefact described above.

The JSON parsing fragility of the AI generation pipeline would benefit
from the adoption of structured output modes now available in the
OpenAI API \cite{openai2024structured}, which constrain the model's
output to a specified JSON schema at the API level rather than relying
on prompt instructions alone. This would eliminate the most common
class of parsing errors encountered during testing. Similar features
are emerging across other LLM providers and would allow the dual
backend architecture to be retained while reducing parser complexity
on both sides.

From an infrastructure perspective, deploying the application to a
cloud provider using the existing Docker Compose configuration as a
starting point, and adding HTTPS termination, a managed database, and
horizontal scaling for the backend, would make the platform suitable
for real user traffic. The continuous integration pipeline built
during the project provides a strong foundation for automating this
deployment process. The code execution engine would also need to be
replaced with a properly sandboxed runner such as a dedicated judging
system or an isolated container per submission before the platform
could responsibly accept code from untrusted users.

A further direction, tangential but interesting, would be to feed the
enriched annotations back into the generation prompt. Once the system
knows that it has just produced an Array problem of medium difficulty,
the next request could be steered explicitly towards a different
topic or a harder level, producing a more balanced practice session
without requiring the user to manage that breadth manually. This
closes the loop between the generative and discriminative components
of the system in a way that the current architecture hints at but
does not yet exploit.

\section{Closing Remarks}

This project demonstrates that it is entirely feasible to build a full
stack educational platform in which AI generated content is
automatically enriched by custom trained machine learning models, all
within the scope of a final year applied project. The combination of
large language model generation with lightweight, fast scikit learn
classifiers \cite{pedregosa2011} proves to be a practical and
effective architecture: the LLM handles the creative and domain
general task of producing coherent questions, while the classifiers
handle the structured annotation task of labelling them.

The limitations identified throughout this dissertation, the proxy
label ceiling, the class imbalance in the topic tagger, the JSON
parsing edge cases, are not failures of ambition but markers of where
the next iteration of this work begins. This project delivers a fully realised platform that integrates generation, tagging, grading, and execution of programming challenges with a social and gamified learning experience.